\section{总结与延伸}
Graham 的这场访谈可以归结为三个平行推进的议题：架构、评测与治理。每条道路都需要大量工程实践、全面评估与安全反思，才能让 Agent 在高风险的软件开发中被广泛采用。

\begin{table}[H]
\centering
\begin{tabular}{p{0.34\textwidth} p{0.58\textwidth}}
\toprule
主题 & 要点 \\
\midrule
背景与动机 & 软件正在吞噬各行业，把“写代码”的能力下放给 Agent 能让更多团队带着工程思维去解决问题。 \\
\midrule
架构与评测 & OpenHands 提供透明 Agent 平台，SWE-bench Lite/Full 提供逐步递进的工程评估指标。 \\
\midrule
未来与治理 & 需要 agent 风格的数据、更多评测维度与内建的安全机制，才能让 Agent 安全进入主流开发流程。 \\
\bottomrule
\end{tabular}
\caption{本次演讲的三个核心要素。}
\end{table}

进一步阅读：
\begin{itemize}
    \item OpenHands 开源仓库：\href{https://github.com/OpenHands/OpenHands}{https://github.com/OpenHands/OpenHands}。
    \item SWE-bench 论文：Jimenez et al., ``SWE-bench: Can Language Models Resolve Real-World GitHub Issues?'' ICLR 2024。
    \item All Hands AI 官方博客与社区反馈入口，关注 agent 安全与人机协作议题。
\end{itemize}

\subsection{本章小结}
\begin{itemize}
    \item 演讲强调的架构、评测与治理三条线必须并行推进，才能把 Agent 安全推向生产。
    \item 关注 OpenHands 与 SWE-bench 的演化，有助于把代理式软件开发落地。
\end{itemize}

\end{document}
